% =====================================================================
\chapter{Proposed Work}
\label{ch:proposed}
% =====================================================================

\section{Introduction}

This chapter presents the approach taken by \veriscript{} and justifies each
design decision with reference to the literature reviewed in
Chapter~\ref{ch:lit}. It then gives the methodology as a sequence of steps,
states the training, inference and threshold-tuning procedures as algorithms,
and describes each step in detail.

\subsection{Overview of the Proposed Approach}

\veriscript{} is a two-stage detector. Every document is first scored by a
\emph{heuristic stage}: a small classifier over perplexity, burstiness and
stylometric features whose output is calibrated into a probability $p_h(d)$.
A \emph{confidence gate} then decides whether that probability is reliable
enough to act on. If $p_h(d)$ is clearly low or clearly high, the decision is
returned immediately together with the features that drove it. If $p_h(d)$
lies inside the ambiguity band $[\Tlow,\Thigh]$, the document is passed to a
\emph{transformer stage}, a fine-tuned RoBERTa-base (or DeBERTa-v3) classifier,
whose probability $p_t(d)$ becomes the final answer. The band is not chosen
by hand; it is tuned on a validation set to meet the accuracy and fairness
constraints of Equation~\eqref{eq:optproblem} at the lowest expected cost.

\subsection{Justification of the Design}

\textbf{Why a cascade and not a single model?} The review of efficient
inference showed that, when the difficulty of inputs varies, routing by
confidence preserves most of an expensive model's accuracy while avoiding most
of its cost~\cite{viola2001rapid,teerapittayanon2016branchynet,xin2020deebert,
chen2023frugalgpt}. AI-text detection has exactly this property. Unedited
output from a chat assistant is statistically very different from human text,
while edited, paraphrased or very formal texts are hard. A single heuristic
model would be wrong too often on the hard cases; a single transformer would
waste computation on the easy ones.

\textbf{Why a cascade and not an ensemble?} An alternative way of combining
the two families is to run both on every document and average or stack their
outputs. Such an ensemble may be slightly more accurate, but it always pays
the full cost of the transformer, so it does not address problem P3. It also
blurs the explanation, because every decision depends on the opaque model.
The cascade preserves a clean separation: fast-path decisions are fully
explainable by construction.

\textbf{Why these heuristic features?} Perplexity and token-rank statistics are
the most consistently useful zero-cost signals in the literature, from
GLTR~\cite{gehrmann2019gltr} to Binoculars~\cite{hans2024binoculars}.
Burstiness captures the uniformity of generated text that perplexity alone
misses. Stylometric features are free to compute and were shown on HC3 to
separate ChatGPT from human answers~\cite{guo2023hc3}. Learned combinations of
such features, as in Ghostbuster~\cite{verma2024ghostbuster}, generalise better
than any single statistic. The heuristic stage therefore uses one small
reference language model (GPT-2 small, about 124 million parameters) to obtain
all probability-based features in a single forward pass, and adds features that
need no model at all. It deliberately avoids methods that need many model calls
(DetectGPT~\cite{mitchell2023detectgpt}) or two large models
(Binoculars), because the first stage must be cheap.

\textbf{Why a linear or tree-based head?} Logistic regression gives a weight to
every feature, so the contribution of each feature to a decision can be read
off directly; gradient-boosted trees~\cite{chen2016xgboost} capture
interactions between features (for example, low perplexity is more suspicious
when burstiness is also low) at little extra cost. Both are trained in
seconds and run in microseconds.

\textbf{Why calibrate?} The gate compares $p_h(d)$ with fixed thresholds.
Uncalibrated scores, especially from boosted trees, are systematically
over- or under-confident~\cite{guo2017calibration}, which would make the
thresholds meaningless and the escalation rate unpredictable. Platt
scaling~\cite{platt1999probabilistic} or isotonic regression on a held-out
split corrects this.

\textbf{Why RoBERTa-base and DeBERTa-v3?} Fine-tuned RoBERTa has been the
standard strong detector since the GPT-2 release~\cite{solaiman2019release} and
remains a strong baseline on M4, MAGE and RAID~\cite{wang2024m4,li2024mage,
dugan2024raid}. DeBERTa-v3~\cite{he2023debertav3} is generally stronger on
classification tasks at similar size and is evaluated as an alternative. Both
are base-sized models that can run on a CPU, which keeps the escalation path
usable on modest institutional hardware. Larger models (RoBERTa-large as used by
RADAR~\cite{hu2023radar}) are kept as an optional upgrade.

\textbf{How does the design address fairness?} Liang \textit{et al.}
showed that low perplexity and low lexical diversity drive false positives on
non-native writing~\cite{liang2023biased}. In the cascade, a human essay that
looks moderately predictable falls into the ambiguity band rather than being
labelled AI-generated by the heuristic stage, because the band is tuned so that
the fast path is used only where the heuristic is reliably correct. The
fairness constraint in Equation~\eqref{eq:optproblem} makes this explicit: a
band that achieves low cost by passing biased fast-path decisions is rejected
during tuning.

\subsection{Design Principles}

The design follows five principles derived from the objectives:
\begin{enumerate}
  \item \textbf{Cheap by default}: every document pays only for the heuristic
        stage unless there is a reason to pay more.
  \item \textbf{Explain what you decide}: any decision made without the
        transformer is accompanied by the feature values that produced it.
  \item \textbf{Measure, do not assume}: thresholds, escalation rate,
        generalisation and fairness are all measured on held-out data.
  \item \textbf{Modularity}: the heuristic stage, the gate and the transformer
        stage are independent components with fixed interfaces, so that either
        model can be replaced (for example by Fast-DetectGPT features or by a
        RADAR-style encoder) without changing the rest of the system.
  \item \textbf{Local and open}: all components run locally on openly licensed
        software; no document leaves the institution.
\end{enumerate}

% ---------------------------------------------------------------------
\section{Proposed Methodology}
% ---------------------------------------------------------------------

The methodology is organised into three stages that correspond to objectives
O1--O3, followed by evaluation (O4) and packaging (O5), as shown in
Fig.~\ref{fig:methodology}. The steps are:

\begin{description}[leftmargin=1.6cm,style=nextline,itemsep=3pt]
  \item[Step 1:] \textbf{Corpus construction.} Assemble a balanced human/AI
        corpus from HC3, M4 and RAID (with optional custom GPT-4, Claude and
        Llama generations); stratify by domain and generator; create training,
        validation, calibration and test splits, including a held-out
        generator split.
  \item[Step 2:] \textbf{Pre-processing.} Normalise text, segment sentences,
        filter short or duplicated samples, and remove obvious artefacts
        (for example ``As an AI language model'') that would make the task
        trivially easy.
  \item[Step 3:] \textbf{Heuristic feature extraction.} Compute perplexity,
        token-rank, burstiness and stylometric features for every document.
  \item[Step 4:] \textbf{Heuristic classifier training and calibration.} Train
        logistic regression and gradient-boosted trees; calibrate the better
        model on the calibration split to obtain $p_h(d)$.
  \item[Step 5:] \textbf{Transformer fine-tuning.} Fine-tune RoBERTa-base
        (DeBERTa-v3 is supported as an alternative) on the training split; keep
        the checkpoint with the best validation F1; record latency.
  \item[Step 6:] \textbf{Threshold tuning.} Search for the ambiguity band
        $[\Tlow,\Thigh]$ that minimises expected cost subject to the accuracy
        and fairness constraints on the validation split.
  \item[Step 7:] \textbf{Cascade inference.} At run time, apply the heuristic
        stage, the gate and, where needed, the transformer stage; produce the
        label, probability, route and explanation.
  \item[Step 8:] \textbf{Evaluation and packaging.} Evaluate the three systems
        on the test splits; package the cascade as a CLI and web application.
\end{description}

\begin{figure}[H]
\centering
\begin{tikzpicture}[node distance=0.45cm and 0.5cm]
  \node[blk,text width=3.4cm] (data) {\textbf{Step 1--2}\\Corpus construction and pre-processing};
  \node[blk,text width=3.1cm,right=1.0cm of data,yshift=1.3cm] (s1) {\textbf{Stage I (O1)}\\Steps 3--4: heuristic features, head, calibration};
  \node[blk,text width=3.1cm,right=1.0cm of data,yshift=-1.3cm] (s2) {\textbf{Stage II (O2)}\\Step 5: transformer fine-tuning};
  \node[blk,text width=3.0cm,right=1.0cm of s1,yshift=-1.3cm] (s3) {\textbf{Stage III (O3)}\\Steps 6--7: gate and threshold tuning};
  \node[blk,text width=3.4cm,below=1.6cm of s3] (ev) {\textbf{Step 8 (O4)}\\Unified evaluation: accuracy, efficiency, generalisation, fairness};
  \node[blk,text width=3.4cm,left=1.0cm of ev] (pk) {\textbf{Step 8 (O5)}\\Packaging: CLI and web demo};
  \draw[arr] (data.east) -- ++(0.4,0) |- (s1.west);
  \draw[arr] (data.east) -- ++(0.4,0) |- (s2.west);
  \draw[arr] (s1.east) -- ++(0.3,0) |- (s3.west);
  \draw[arr] (s2.east) -- ++(0.3,0) |- (s3.west);
  \draw[arr] (s3) -- (ev);
  \draw[arr] (ev) -- (pk);
\end{tikzpicture}
\caption{Overall methodology of the proposed work}
\label{fig:methodology}
\end{figure}

\subsection{Training Procedure}

Algorithm~\ref{alg:train} summarises how the two stages are trained. The
corpus is split so that the calibration and validation sets are disjoint from
the training set; otherwise the calibration would be fitted to over-confident
in-sample predictions.

\begin{algorithm}[H]
\caption{Training the VeriScript components}
\label{alg:train}
\setstretch{1.15}
\begin{algorithmic}[1]
\Require Labelled corpus $\mathcal{D}$; reference LM $\theta_{\mathrm{ref}}$; pretrained encoder $\theta_{\mathrm{enc}}$
\Ensure Calibrated heuristic model $g_h$; fine-tuned transformer $g_t$
\State Split $\mathcal{D}$ into $\mathcal{D}_{tr}$, $\mathcal{D}_{cal}$, $\mathcal{D}_{val}$, $\mathcal{D}_{te}$ (stratified by generator and domain)
\ForAll{$d \in \mathcal{D}$}
   \State $\mathbf{h}(d) \gets \Call{ExtractFeatures}{d,\theta_{\mathrm{ref}}}$ \Comment{one LM forward pass}
\EndFor
\State Standardise features using statistics of $\mathcal{D}_{tr}$
\State Train LR and GBDT heads on $\{(\mathbf{h}(d),y)\}_{d\in\mathcal{D}_{tr}}$
\State Select the head with higher AUROC on $\mathcal{D}_{val}$
\State Fit calibrator (Platt or isotonic) on $\mathcal{D}_{cal}$; $g_h \gets$ calibrated head
\State Initialise $g_t$ from $\theta_{\mathrm{enc}}$ with a linear classification layer
\For{epoch $= 1$ \textbf{to} $E$}
   \ForAll{mini-batch $B \subset \mathcal{D}_{tr}$}
      \State Minimise cross-entropy of $g_t$ on $B$ with AdamW
   \EndFor
   \State Evaluate F1 on $\mathcal{D}_{val}$; keep best checkpoint
\EndFor
\State \Return $g_h$, $g_t$
\end{algorithmic}
\end{algorithm}

\subsection{Cascade Inference Procedure}

Algorithm~\ref{alg:infer} gives the inference procedure, which implements
Equation~\eqref{eq:cascade}, and Fig.~\ref{fig:flowchart} shows the same
logic as a flowchart.

\begin{algorithm}[H]
\caption{Confidence-gated cascade inference}
\label{alg:infer}
\setstretch{1.15}
\begin{algorithmic}[1]
\Require Document $d$; models $g_h$, $g_t$; band $[\Tlow,\Thigh]$; threshold $\tau$
\Ensure Label $\hat{y}$, probability $p$, route $r$, explanation $e$
\State $d \gets \Call{Preprocess}{d}$
\If{$\Call{WordCount}{d} < 100$}
   \State \Return (\textsc{Undetermined}, --, --, ``insufficient text'')
\EndIf
\State $\mathbf{h} \gets \Call{ExtractFeatures}{d}$;\quad $p_h \gets g_h(\mathbf{h})$
\If{$p_h < \Tlow$ \textbf{or} $p_h > \Thigh$}
   \State $p \gets p_h$;\quad $r \gets$ \texttt{fast\_path}
   \State $e \gets \Call{TopFeatures}{g_h,\mathbf{h},k=5}$
\Else
   \State $W \gets \Call{SlidingWindows}{d, L=512, \text{stride}=384}$
   \State $p \gets \frac{1}{|W|}\sum_{w\in W} g_t(w)$;\quad $r \gets$ \texttt{escalated}
   \State $e \gets \{\,p_h,\ p,\ \Call{TopFeatures}{g_h,\mathbf{h},k=5}\,\}$
\EndIf
\State $\hat{y} \gets \mathbb{1}[p \ge \tau]$
\State \Return $(\hat{y}, p, r, e)$
\end{algorithmic}
\end{algorithm}

\begin{figure}[H]
\centering
\begin{tikzpicture}[node distance=0.55cm and 0.9cm]
  \node[term] (start) {Start};
  \node[io,below=of start] (in) {Read document $d$};
  \node[blk,below=of in,text width=4.2cm] (pre) {Pre-process: normalise, split sentences};
  \node[dec,below=of pre] (len) {Words $\ge$ 100?};
  \node[term,right=1.2cm of len] (short) {Report ``insufficient text''};
  \node[blk,below=of len,text width=4.2cm] (feat) {Extract features $\mathbf{h}(d)$; compute $p_h(d)$};
  \node[dec,below=of feat] (gate) {$\Tlow \le p_h \le \Thigh$?};
  \node[blk,left=1.0cm of gate,text width=3.1cm,yshift=-1.6cm] (fast) {Fast path:\\$p \gets p_h$;\\explain top features};
  \node[blk,right=1.0cm of gate,text width=3.1cm,yshift=-1.6cm] (esc) {Escalate:\\windowed transformer;\\$p \gets p_t$};
  \node[blk,below=2.5cm of gate,text width=4.2cm] (dec) {$\hat{y} \gets \mathbb{1}[p \ge \tau]$; log route and latency};
  \node[io,below=of dec] (out) {Output label, $p$, route, explanation};
  \node[term,below=of out] (stop) {Stop};
  \draw[arr] (start) -- (in);
  \draw[arr] (in) -- (pre);
  \draw[arr] (pre) -- (len);
  \draw[arr] (len) -- node[lbl,above] {No} (short);
  \draw[arr] (len) -- node[lbl,right] {Yes} (feat);
  \draw[arr] (feat) -- (gate);
  \draw[arr] (gate.west) -| node[lbl,above,pos=0.25] {No} (fast.north);
  \draw[arr] (gate.east) -| node[lbl,above,pos=0.25] {Yes} (esc.north);
  \draw[arr] (fast.south) |- (dec.west);
  \draw[arr] (esc.south) |- (dec.east);
  \draw[arr] (dec) -- (out);
  \draw[arr] (out) -- (stop);
\end{tikzpicture}
\caption{Flowchart of the VeriScript cascade inference algorithm}
\label{fig:flowchart}
\end{figure}

\subsection{Threshold Tuning Procedure}

The ambiguity band is selected by a grid search on the validation set, as in
Algorithm~\ref{alg:tune}. For each candidate band the escalation rate,
cascade F1 and fairness disparity are computed using the already-stored
predictions of both stages, so the search requires no further model inference
and finishes in seconds.

\begin{algorithm}[H]
\caption{Validation-based tuning of the ambiguity band}
\label{alg:tune}
\setstretch{1.15}
\begin{algorithmic}[1]
\Require Validation predictions $\{(p_h(d), p_t(d), y, g)\}$ where $g$ marks native/non-native; tolerances $\varepsilon$, $\delta$; grid step $s$
\Ensure Band $(\Tlow^{*},\Thigh^{*})$
\State $F_1^{t} \gets$ F1 of the transformer-only detector on the validation set
\State $best \gets (0,1)$;\quad $\alpha_{best} \gets 1$
\For{$\Tlow \gets 0$ \textbf{to} $0.5$ \textbf{step} $s$}
  \For{$\Thigh \gets 0.5$ \textbf{to} $1$ \textbf{step} $s$}
     \State Form cascade predictions by Eq.~\eqref{eq:cascade}
     \State Compute $\alpha$, $F_1^{c}$ and $\Delta_{FPR}$
     \If{$F_1^{c} \ge F_1^{t}-\varepsilon$ \textbf{and} $\Delta_{FPR} \le \delta$ \textbf{and} $\alpha < \alpha_{best}$}
        \State $best \gets (\Tlow,\Thigh)$;\quad $\alpha_{best} \gets \alpha$
     \EndIf
  \EndFor
\EndFor
\State \Return $best$
\end{algorithmic}
\end{algorithm}

\textbf{Form of the fairness constraint.} Algorithm~\ref{alg:tune} is written
with the disparity constraint $\Delta_{FPR} \le \delta$ of
Eq.~\eqref{eq:optproblem}. A disparity bound alone has a weakness: it can be
satisfied by making the detector \emph{worse} for the better-served group,
for example by raising the false positive rate on native writing until it
matches the rate on non-native writing. The implementation therefore offers
two forms of the constraint, evaluated on a fairness validation set of native
and non-native human essays that is disjoint from the fairness test set:
\begin{itemize}
  \item \emph{disparity}: $\Delta_{FPR} \le \delta$, with $\delta$ equal to the
        transformer-only disparity plus one standard error of that estimate
        (the fairness sets are small, so the reference disparity is itself
        uncertain), and additionally
        $\mathrm{FPR}^{c}_{\text{non-native}} \le \mathrm{FPR}^{t}_{\text{non-native}}$;
  \item \emph{worst-group} (default): the cascade must not flag non-native
        writers more often than the transformer-only detector,
        $\mathrm{FPR}^{c}_{\text{non-native}} \le \mathrm{FPR}^{t}_{\text{non-native}}$,
        which is the second half of hypothesis H4 and cannot be met by harming
        the native group.
\end{itemize}
Among bands that satisfy the constraints, ties in escalation rate are broken
in favour of the higher validation F1. Chapter~\ref{ch:results} explains,
with measured numbers, why the worst-group form was adopted for deployment
and reports the bands selected by both forms.

% ---------------------------------------------------------------------
\section{Description of Each Step}
% ---------------------------------------------------------------------

\subsection{Corpus Construction and Pre-processing (Steps 1--2)}

The corpus combines public benchmarks (and, optionally, custom generations),
as detailed in Chapter~6. Three principles guide its construction. First,
\emph{balance}: human and AI classes are of equal size within each domain and
length bucket, so that the classifier cannot exploit class priors that differ
by topic or length. Second, \emph{diversity}: AI text comes from several
generator families (the GPT family including ChatGPT and GPT-4, Cohere, Dolly,
BLOOMZ, Flan-T5, Llama, MPT and Mistral) and several decoding settings,
because detectors trained on one generator transfer
poorly~\cite{dugan2024raid,wang2024m4}. Third, \emph{held-out evaluation}: one
generator family is removed entirely from training and calibration and is used
only for testing, which implements the leave-one-generator-out protocol.

Pre-processing applies Unicode normalisation (NFKC), collapses whitespace,
strips HTML or Markdown markup and removes boilerplate phrases that reveal the
source trivially, such as refusal templates. Documents shorter than 100 words
are excluded from training, and near-duplicates are removed by hashing
normalised text so that the same answer cannot appear in both training and
test sets. Long documents are kept intact for feature extraction; for the
transformer they are split into windows at inference time.

\subsection{Heuristic Feature Extraction (Step 3)}

Each document is passed once through the reference model (GPT-2 small) to
obtain the log-probability, rank and predictive entropy of every token. From
these and from the raw text, the features in Table~\ref{tab:features} are
computed.

\begin{table}[H]
\caption{Heuristic features computed for each document}
\label{tab:features}
\centering
\small
\begin{tabularx}{\textwidth}{|p{0.6cm}|p{3.6cm}|L|p{2.1cm}|}
\hline
\textbf{No.} & \textbf{Feature} & \textbf{Definition} & \textbf{Group} \\ \hline
1 & Mean log-probability & $\frac{1}{N}\sum_i \log P(x_i\mid x_{<i})$ & Perplexity \\ \hline
2 & Document perplexity & Eq.~\eqref{eq:perplexity} & Perplexity \\ \hline
3 & Std.\ of token log-prob & Spread of per-token log-probabilities & Perplexity \\ \hline
4 & Min.\ sentence perplexity & Lowest perplexity over sentences & Perplexity \\ \hline
5 & Max.\ sentence perplexity & Highest perplexity over sentences & Perplexity \\ \hline
6 & Top-10 rank fraction & Share of tokens ranked in the top 10 & Rank (GLTR) \\ \hline
7 & Top-100 rank fraction & Share of tokens ranked 11--100 & Rank (GLTR) \\ \hline
8 & Tail rank fraction & Share of tokens ranked above 1000 & Rank (GLTR) \\ \hline
9 & Mean predictive entropy & Average of $H_i$ (Eq.~\eqref{eq:entropy}) & Entropy \\ \hline
10 & Sentence-length burstiness & $\beta$ of Eq.~\eqref{eq:burstiness} & Burstiness \\ \hline
11 & Perplexity burstiness & Coefficient of variation of sentence perplexities & Burstiness \\ \hline
12 & Type-token ratio & Eq.~\eqref{eq:ttr}, on the first 200 words & Stylometric \\ \hline
13 & Mean word length & Average characters per word & Stylometric \\ \hline
14 & Function-word rate & Share of tokens in a function-word list & Stylometric \\ \hline
15 & Punctuation profile & Rates of commas, semicolons, colons, dashes per sentence & Stylometric \\ \hline
16 & POS bigram profile & Normalised counts of the 20 most informative POS bigrams & Stylometric \\ \hline
\end{tabularx}
\end{table}

Two details matter for robustness. The type-token ratio is computed on a
fixed-length prefix because TTR falls with document length; computing it on
the whole document would make it a proxy for length. Sentence-level perplexity
statistics are computed only over sentences of at least five tokens to avoid
noise from headings and fragments. All features are standardised using the
mean and standard deviation of the training split:
\begin{equation}
  \tilde{h}_j(d) = \frac{h_j(d) - \mu_j}{\sigma_j}.
  \label{eq:standardise}
\end{equation}

\subsection{Heuristic Classifier and Calibration (Step 4)}

\textbf{Logistic regression.} The baseline head is an $\ell_2$-regularised
logistic regression,
\begin{equation}
  s_h(d) = \sigma\!\left(\mathbf{w}^{\top}\tilde{\mathbf{h}}(d) + b\right),
  \label{eq:lr}
\end{equation}
trained by minimising the regularised log-loss
\begin{equation}
  \mathcal{L}(\mathbf{w},b) = -\frac{1}{n}\sum_{k=1}^{n}\Big[y_k \log s_h(d_k) + (1-y_k)\log\big(1-s_h(d_k)\big)\Big] + \frac{\lambda_2}{2}\lVert\mathbf{w}\rVert_2^2 .
  \label{eq:logloss}
\end{equation}

\textbf{Gradient-boosted trees.} The second head is an ensemble of $M$
regression trees built stagewise~\cite{chen2016xgboost}:
\begin{equation}
  F_M(\tilde{\mathbf{h}}) = \sum_{m=1}^{M} \eta\, f_m(\tilde{\mathbf{h}}), \qquad
  s_h(d) = \sigma\!\big(F_M(\tilde{\mathbf{h}}(d))\big),
  \label{eq:gbdt}
\end{equation}
where each tree $f_m$ is fitted to the gradient of the log-loss of the current
ensemble and $\eta$ is the learning rate. Tree depth, number of trees and
learning rate are chosen by cross-validation on the training split.

\textbf{Calibration.} The selected head's score is calibrated on the
calibration split with Platt scaling (Eq.~\eqref{eq:platt}), or with isotonic
regression if the calibration split is large enough. Calibration quality is
checked with the expected calibration error,
\begin{equation}
  \mathrm{ECE} = \sum_{b=1}^{B} \frac{|S_b|}{n}\,\big|\,\mathrm{acc}(S_b) - \mathrm{conf}(S_b)\,\big|,
  \label{eq:ece}
\end{equation}
where $S_b$ is the set of documents whose predicted probability falls in the
$b$-th of $B$ equal-width bins~\cite{guo2017calibration}. The calibrated output
is $p_h(d)$.

\subsection{Transformer Fine-Tuning (Step 5)}

RoBERTa-base is fine-tuned with a linear classification layer on the pooled
first-token representation (Eq.~\eqref{eq:clshead}), using the binary
cross-entropy loss and the AdamW optimiser~\cite{loshchilov2019adamw} with a
linear learning-rate warm-up followed by linear decay. Training documents
longer than the maximum sequence length $L=512$ are truncated to their first
$L$ tokens during training. At inference, a long document is split into
overlapping windows $w_1,\ldots,w_K$ of length $L$ with a stride of 384 tokens,
and the window probabilities are averaged:
\begin{equation}
  p_t(d) = \frac{1}{K}\sum_{k=1}^{K} \sigma\!\left(\mathbf{w}_t^{\top}\mathbf{h}_{\mathrm{CLS}}(w_k) + b_t\right).
  \label{eq:windows}
\end{equation}
The same procedure applies unchanged to DeBERTa-v3-base, which the
implementation supports through a single command-line option; in the
experiments reported here only RoBERTa-base was trained (Chapter~6 gives the
reason). The hyperparameters used are listed in Chapter~6.

\subsection{Confidence Gate and Threshold Tuning (Step 6)}

The gate divides the probability axis into three regions, shown in
Fig.~\ref{fig:band}. Documents whose $p_h(d)$ falls below $\Tlow$ are labelled
human-written on the fast path; those above $\Thigh$ are labelled AI-generated
on the fast path; those between are escalated. The band need not be symmetric
around 0.5. Because a false accusation is more costly than a missed detection
in the academic setting, the tuning procedure typically places $\Thigh$ further
from 0.5 than $\Tlow$, so that the heuristic is allowed to label a document
as AI-generated only when it is very confident.

\begin{figure}[H]
\centering
\begin{tikzpicture}[x=12cm,y=1cm]
  \fill[pattern=north east lines] (0.30,0) rectangle (0.82,0.6);
  \draw[thick] (0,0) -- (1,0);
  \draw[thick] (0,0.6) -- (1,0.6);
  \draw[thick] (0,0) -- (0,0.6) (1,0) -- (1,0.6);
  \draw[very thick] (0.30,-0.2) -- (0.30,0.8);
  \draw[very thick] (0.82,-0.2) -- (0.82,0.8);
  \node[below] at (0,-0.05) {0};
  \node[below] at (1,-0.05) {1};
  \node[below] at (0.30,-0.2) {$\Tlow$};
  \node[below] at (0.82,-0.2) {$\Thigh$};
  \node[below] at (0.56,-0.85) {calibrated heuristic probability $p_h(d)$};
  \node[above,font=\small,align=center] at (0.15,0.75) {Fast path:\\human-written};
  \node[above,font=\small,align=center] at (0.56,0.75) {Escalate to transformer\\(ambiguity band)};
  \node[above,font=\small,align=center] at (0.91,0.75) {Fast path:\\AI-generated};
\end{tikzpicture}
\caption{Partition of the heuristic probability axis by the confidence gate (illustrative band)}
\label{fig:band}
\end{figure}

An equivalent unconstrained form of the tuning objective, useful for plotting
the trade-off curve, is the Lagrangian
\begin{equation}
  J(\Tlow,\Thigh) = \big(1 - F_1^{\text{cascade}}\big) + \lambda\,\frac{\mathbb{E}[C]}{C_h + C_t},
  \label{eq:lagrangian}
\end{equation}
where $\lambda \ge 0$ sets how much accuracy the institution is willing to
trade for a unit of relative cost. Sweeping $\lambda$ traces a curve of
operating points from the heuristic-only system ($\lambda$ large) to the
transformer-only system ($\lambda = 0$). Instead of one fixed setting,
\veriscript{} can therefore offer a small menu of operating points, for
example ``fast'', ``balanced'' and ``thorough'', each with its measured
escalation rate, accuracy and latency.

\subsection{Explanation Generation (Step 7)}

For fast-path decisions made by the logistic-regression head, the contribution
of feature $j$ to the log-odds is exactly $c_j = w_j\,\tilde{h}_j(d)$, so the
explanation lists the $k$ features with the largest $|c_j|$ together with their
raw values and their typical ranges for human and AI text in the training
data, for example: ``perplexity 11.2 (human texts in training: median 34)''.
For the gradient-boosted head, per-feature contributions are obtained from the
tree structure using the path-based attribution built into the XGBoost
library. For escalated decisions the explanation reports the transformer
probability together with the heuristic context, and states clearly that the
decision was made by the neural model. In all cases the output includes a
reminder that the result is a probability and should be reviewed by a person.

\subsection{Evaluation and Packaging (Step 8)}

The three systems (heuristic-only, transformer-only and cascade) are evaluated
on identical test splits for accuracy (Precision, Recall, F1, AUROC),
efficiency (mean and p95 latency, throughput, model size), cross-generator
generalisation (leave-one-generator-out) and fairness (false-positive-rate
disparity on non-native writing). The principal figure of merit is the
fraction of the transformer's accuracy advantage over the heuristic that the
cascade recovers, per unit of latency saved:
\begin{equation}
  \text{Recovery} = \frac{F_1^{\text{cascade}} - F_1^{\text{heuristic}}}{F_1^{\text{transformer}} - F_1^{\text{heuristic}}},
  \qquad
  \text{Speed-up} = \frac{\bar{t}_{\text{transformer}}}{\bar{t}_{\text{cascade}}}.
  \label{eq:recovery}
\end{equation}
The packaged system exposes the cascade through a command line tool for batch
processing and a web interface for single documents, as described in
Chapters~\ref{ch:design} and~6.

\section{Summary}

This chapter presented the proposed confidence-gated cascade and justified its
components using the literature: a single-pass heuristic stage over sixteen
interpretable features, calibrated so that its probability can be thresholded;
a fine-tuned RoBERTa-base or DeBERTa-v3 classifier for ambiguous documents;
and a validation-based search for the ambiguity band under accuracy and
fairness constraints. The training, inference and tuning procedures were given
as algorithms, and each step was described. The next chapter specifies the
system that implements this methodology.
